\begin{frame}
\frametitle{Skryté paměti -- cache}

\begin{itemize}
\item je označení pro vyrovnávací paměť používanou ve výpočetní technice
\item Zařazujeme ji mezi dva subsystémy s různou rychlostí. Vyrovnává se jí rychlost přístupu k informacím.
\item Účelem skryté paměti je urychlit přístup k často používaným datům na „pomalých“ médiích jejich překopírováním na média rychlá.
\item Většinou z velké části pracuje automaticky a přispůsobuje se aktuální potřebě programu.
\item Přesto je nutné o její existenci vědět a často potřebuje servisní obsluhu na úrovni operačního systému a v něktěrých případech i programů.
\item Při nepromyšlené volbě datových struktur a algoritmů se její efekt ztrácí.
\end{itemize}

\end{frame}

\begin{frame}
\frametitle{Skryté paměti -- terminologie}

\begin{itemize}
\item \textbf{Cache hit} (zásah) pojmenování situace, kdy požadovaná hodnota ve skryté paměti (cache) je.
\item \textbf{Cache miss}, opak, \textbf{výpadek}, data v cache ještě nejsou.
\item \textbf{Cache line} (řádka) nebo \textbf{Cache block} – základní kopírovatelná jednotka mezi hierarchickými úrovněmi. 
\item V praxi se velikost \textbf{Cache} Line pohybuje od 8B do 1KB, typicky 64B.
\item \textbf{Hit rate} -- počet paměťových přístupů obsloužených danou úrovní cache dělený všemi přístupy
\item \textbf{Miss rate} -- poměr přístupů, které je potřeba obsloužit z pomalejší paměti = 1 - Hit rate
\item \textbf{Average Memory Access Time} (AMAT) $$AMAT = Hit Time + Miss Rate \times Miss Penalty$$
\item AMAT pro víceúrovňovou cache lze spočítat rekurzivní aplikací výšeuvedeného vztahu
\end{itemize}

\end{frame}

\begin{frame}
\frametitle{Skryté paměti -- realizace}
{
\centering

\includegraphics[width=0.70\linewidth]{generated/memory/cache/cache-principle.pdf}

}
\vskip 2mm

\begin{itemize}
\item \textbf{Tag} je index odpovídajícího bloku v operační paměti (v podstatě se jedná o hodnotu ukazatele/adresy dělenou délkou bloku).
\item \textbf{Data} pole obsahující vlastní hodnoty na příslušné/ných adrese/ách.
\item \textbf{Validity bit} – bit platnosti. Indikuje, zda je obsah pole Data vůbec platný.  
\item \textbf{Dirty bit} – rozšiřující pole v obsahu paměti. Indikuje, že v cache (cache) je jiná hodnota, než v paměti hlavní.
\end{itemize}

\end{frame}


\begin{frame}
\frametitle{Zpracování výpadku skryté paměti, cache miss}

\begin{itemize}
\item Data musí být nahrána z hlavní paměti, obvykle jsou již ale všechny položky vyrovnávací paměťi (cache) zaplněné daty z přetrvávajícími předchozího běhu programu.
\item Některou z položek, kterou je možné využít pro uložení dat z dané adresy je potřeba uvolnit.
\item Na výběru položky k vyřazení velmi záleží, pokud bude vybraná ta, která bude opětovně potřeba, dojde k snížení výkonu.
\item \textbf{Cache replacement policy} -- pravidla pro výběr položky k vyřazení
\begin{itemize}
\item \textbf{Random} -- vybraná je náhodná položka
\item \textbf{LRU} (Least Recently Used) -- vybrána je nejdelší dobu nepoužitá položka, do obvodů cache musí být ke každé skupině bloků přidané další informace, které umožňují sledovat pořadí posledních přístupů k jednotlivým položkám.
\item \textbf{LFU} (Least Frequently Used) -- sleduje se, jak často/kolikrát se k položkám přistupuje, vyžaduje přidat zapomínání.
\item \textbf{ARC} (Adaptive Replacement Cache) – kombinace LRU a LFU
\end{itemize}
\end{itemize}

\end{frame}

\begin{frame}
\frametitle{Zápisy procesoru do hlavní paměti}

\begin{itemize}
\item V cestě je cache, která může blok, do kterého je zapisováno, obsahovat.
\item Minimálně z pohledu daného procesoru je nutné zajistit koherenci dat pro daný procesor (často i pro více procesorů -- vlákna) pro přístupy ke každé jednotlivé adrese i pokud existuje více cest přístupu 
\item \textbf{Write through} (zápis, propsání zkrz) cache -- pokud se již data v cache nachází, jsou modifikovaná (ve variantě s automatickou alokací jsou i při výpadku nahraná a pak modifikovaná). Data jsou zároveň odeslaná do hlavní paměti, buď přímo nebo přes zápisovou frontu (write buffer)
\item \textbf{Write back} -- data jsou zapsaná do příslušného bloku cache, pokud není v cache, je blok nejdříve načtený. Blok je označený \textbf{dirty bit}em. Zápis do další úrovně až pokud je potřeba danou položku cache uvolnit pro jiná data (replacement) nebo je vyžádaná sychronizace procesorem, systémem (cache flush).
\end{itemize}

\end{frame}

\begin{frame}
\frametitle{Základní typy organizace paměti cache}

Uvažujeme cache o velikosti 8 bloků a kam může být mapovaný přístup na adresu/blok 12 pro tři varianty organizace skryté paměti

\begin{tabular}{p{0.29\linewidth}p{0.29\linewidth}p{0.29\linewidth}}
\textbf{Plně asociativní} & \textbf{Přímo mapovaná} & \textbf{Dvoucestná} \\
\textbf{Fully associative:} & \textbf{Direct mapped} & \textbf{2-way associative} \\
Adresa 12 může být umístěna libovolně &
Adresa 12 může být umístěna jen do bloku 4 (12 mod 8) &
Adresa 12 může být umístěna do sady 0 (12 mod 4) \\
\includegraphics[width=3cm]{generated/memory/cache/cache-schema-fully.pdf} &
\includegraphics[width=3cm]{generated/memory/cache/cache-schema-direct.pdf} &
\includegraphics[width=3cm]{generated/memory/cache/cache-schema-2way.pdf} \\
\end{tabular}

\end{frame}

\begin{frame}
\frametitle{Přímo mapovaná paměť cache -- mapování}

{
\centering

\includegraphics[width=0.9\linewidth]{generated/memory/cache/cache-direct-mapping-1byte.pdf}

}

\end{frame}

\begin{frame}
\frametitle{Přímo mapovaná paměť cache s blokem rovným slovu}

\begin{itemize}
\item \textbf{Capacity} -- C ... kapacita (zde 1024 v bytech a 256 v slovech)
\item \textbf{Number of sets} -- SN .. počet setů, (zde 256)
\item \textbf{Word size} – WS .. velikost slova (zde 4 byte)
\item \textbf{Block size} – BS .. velikost bloku (zde 1 slovo, 4 byte)
\item \textbf{Number of blocks} -- BN .. počet bloků (zde 256)
\item \textbf{Degree of associativity} -- N . stupeň asociativity, počet cest (zde 1)
\end{itemize}

{
\centering

\includegraphics[width=0.70\linewidth]{generated/memory/cache/cache-direct-simple.pdf}

}

C = 1024 bytů, SN = BN = 256, BS = 1, N = 1

\end{frame}


\begin{frame}
\frametitle{Přímomapovaná paměť cache s většími bloky}

BS = 4 (16 byte, slovo WS = 4 byte), setů SN = 4

set = (adresa div BS) mod SN

{
\centering

\includegraphics[width=0.75\linewidth]{generated/memory/cache/cache-direct-ms.pdf}

}

{\tiny Zdroj: Michal Štepanovský}
\end{frame}

\begin{frame}
\frametitle{Obecná organizace paměti cache s N cestami}

{
\centering

\includegraphics[width=0.95\linewidth]{generated/memory/cache/cache-n-way-ms.pdf}

}

{\tiny Zdroj: Michal Štepanovský}

\end{frame}
